\documentclass[a4paper,14pt]{extarticle}
\usepackage{suai-report}

\suaisetup{
  department   = 41,
  teacher-post = Старший преподаватель,
  teacher      = В. В. Боженко,
  type         = Отчет о лабораторной работе,
  number       = 4,
  title        = Кластеризация,
  course       = Введение в анализ данных,
  group        = 4414,
  student      = Г. В. Моисеев,
  date         = today,
}

\begin{document}

\maketitle
\suaitoc

\suaiintro

Цель работы — изучение алгоритмов и методов кластеризации на практике.

\suaitasks
создать набор данных функцией make\_data, стандартизовать его и выполнить кластеризацию методом k-means для трёх кластеров
визуализировать кластеры и их центры, рассчитать коэффициент силуэта
повторить кластеризацию для двух и четырёх кластеров, сравнить результаты и построить график по методу локтя
загрузить набор данных варианта, провести предварительную обработку и выделить целевую переменную
построить матрицу диаграмм рассеяния с выделением значений целевой переменной и стандартизовать числовые признаки
выполнить кластеризацию методом k-means, подобрать оптимальное число кластеров и описать полученные кластеры
рассчитать метрики качества кластеризации
выполнить кластеризацию иерархическим агломеративным методом, построить дендрограмму и сравнить результаты двух методов

\section{Вариант задания}

Номер в списке группы — 16. В методических указаниях приведено 5 вариантов, поэтому номер взят по кругу и выполнен вариант 1: набор данных medical\_diagnosis.csv с информацией о болезнях. Каждая строка набора описывает одного пациента. Столбцы набора приведены в~\tabref{columns}.

\suaitable[columns]{Столбцы набора данных medical\_diagnosis.csv}
Столбец | Содержание
age | возраст пациента, лет
sex | пол: Male, Female
BMI | индекс массы тела, кг/м²
systolic\_bp | систолическое артериальное давление, мм рт. ст.
cholesterol\_mg\_dl | общий холестерин сыворотки, мг/дл
fasting\_glucose\_mg\_dl | глюкоза крови натощак, мг/дл
smoking\_status | статус курения: Never, Former, Current
family\_history | наследственность по заболеванию: Yes, No
activity\_level | уровень физической активности: Low, Moderate, High
symptom\_score | суммарная выраженность симптомов, от 0 до 10
CRP\_mg\_L | C-реактивный белок, мг/л (маркер воспаления)
WBC\_10e9\_L | лейкоциты, $\times 10^9$/л
ECG\_result | результат ЭКГ: Normal, ST-T abnormality, LVH
Disease | целевой класс: 1 — болеет, 0 — не болеет

В части 1 используются данные, созданные функцией make\_data из методических указаний. Значения её параметров в указаниях не заданы, поэтому взято $n = 100$ объектов в каждой группе и зерно генератора случайных чисел $seed = 16$, равное номеру варианта.

Работа выполнена в Jupyter-ноутбуке ЛР4\_Моисеев\_1.ipynb. Ссылка на ноутбук: \url{https://github.com/Echways/ItDA/tree/main/lab-4}.

\section{Кластеризация сгенерированных данных}

\subsection{Создание набора данных}

Функция make\_data взята из методических указаний без изменений. Она создаёт три группы точек с центрами (3; 3), (6; 9) и (9; 3), добавляет к каждой координате нормальный шум с единичным стандартным отклонением, а затем умножает признак x на 20 и признак y на 0,5. После этого центры групп находятся в точках (60; 1,5), (120; 4,5) и (180; 1,5). В конце строки перемешиваются. Код функции и первые строки полученного набора показаны на~\figref{p1_data}.

\suaiimg[\textwidth]{p1_data}{Создание набора данных функцией make\_data}

Получено 300 объектов с двумя признаками. Описательные статистики набора приведены на~\figref{p1_describe}.

\suaiimg[0.19\textwidth]{p1_describe}{Описательные статистики сгенерированного набора}

Признак x меняется от 10,27 до 218,14 при стандартном отклонении 51,28, признак y — от 0,51 до 5,52 при стандартном отклонении 1,47. Разброс x в 35 раз больше разброса y.

\subsection{Стандартизация данных}

Метод k-means относит объект к ближайшему центру по евклидову расстоянию. Без стандартизации расстояние почти полностью определялось бы признаком x, и различие групп по y алгоритм бы не учитывал. Поэтому оба признака приведены к среднему 0 и стандартному отклонению 1 с помощью StandardScaler (рисунок~\ref{fig:p1_scale}).

\suaiimg[\textwidth]{p1_scale}{Стандартизация сгенерированных данных}

\subsection{Кластеризация методом k-means}

Модель KMeans для трёх кластеров обучена методом fit\_predict, который сразу возвращает метки кластеров (рисунок~\ref{fig:p1_fit}). Параметр n\_init равен 10: алгоритм запускается 10 раз с разными начальными центрами и оставляет разбиение с наименьшей суммой квадратов расстояний до центров. Параметр random\_state фиксирует начальные центры, чтобы результат воспроизводился.

\suaiimg[\textwidth]{p1_fit}{Обучение модели k-means для трёх кластеров и метки кластеров}

Каждый объект получил метку 0, 1 или 2. В каждом кластере ровно по 100 объектов — столько же, сколько в каждой сгенерированной группе. Центры кластеров из атрибута cluster\_centers\_ приведены на~\figref{p1_centers} в стандартизованных координатах и в исходных единицах.

\suaiimg[\textwidth]{p1_centers}{Центры кластеров}

В исходных единицах центры находятся в точках (117,70; 4,47), (60,42; 1,59) и (178,02; 1,46). Это близко к центрам, заложенным в функцию make\_data: (120; 4,5), (60; 1,5) и (180; 1,5).

Кластеры показаны на~\figref{p1_clusters}: объекты каждого кластера окрашены своим цветом, центры отмечены чёрными крестами.

\suaiimg[0.75\textwidth]{p1_clusters}{Кластеры и их центры при трёх кластерах}

Три облака точек отделены друг от друга, и каждое целиком окрашено одним цветом. Центры кластеров находятся в середине своих облаков.

Качество разбиения оценено коэффициентом силуэта. Для одного объекта он вычисляется по~\formref{silhouette}:
\suaieq[silhouette]{s_i = \frac{b_i - a_i}{\max(a_i, b_i)}}
a_i | среднее расстояние от объекта $i$ до остальных объектов его кластера
b_i | среднее расстояние от объекта $i$ до объектов ближайшего чужого кластера

Значение близко к 1, когда объект лежит глубоко в своём кластере, около 0 — когда он на границе двух кластеров, и меньше 0, когда чужой кластер ближе своего. Функция silhouette\_score возвращает среднее по всем объектам (рисунок~\ref{fig:p1_silhouette}).

\suaiimg[\textwidth]{p1_silhouette}{Коэффициент силуэта для трёх кластеров}

Коэффициент силуэта равен 0,720: кластеры плотные и хорошо отделены друг от друга.

\subsection{Сравнение разбиений на два, три и четыре кластера}

Обучение, визуализация и расчёт коэффициента силуэта повторены для двух и четырёх кластеров. Результаты для трёх значений n\_clusters показаны на~\figref{p1_compare_plot}.

\suaiimg[\textwidth]{p1_compare_plot}{Кластеры при n\_clusters, равном 2, 3 и 4}

Функция make\_data перемешивает строки, но сохраняет исходный индекс: строки 0–99 относятся к первой группе, 100–199 — ко второй, 200–299 — к третьей. Поэтому номер сгенерированной группы равен целой части от деления индекса на 100, и с ним можно сравнить кластеры. Для сравнения использована метрика ARI, которая равна 1 при совпадении двух разбиений и близка к 0 при случайном совпадении. Сводные показатели приведены на~\figref{p1_compare}, распределение объектов исходных групп по кластерам — на~\figref{p1_crosstab}.

\suaiimg[0.71\textwidth]{p1_compare}{Показатели разбиений на два, три и четыре кластера}

\suaiimg[0.5\textwidth]{p1_crosstab}{Распределение объектов исходных групп по кластерам}

При двух кластерах коэффициент силуэта равен 0,499. Исходные группы 0 и 1 объединены в один кластер из 200 объектов, и его центр оказался в пустом месте между двумя облаками. При трёх кластерах коэффициент наибольший, 0,720, каждый кластер совпадает с одной исходной группой, ARI~= 1. При четырёх кластерах коэффициент равен 0,585: исходная группа 0 разрезана на кластеры из 52 и 49 объектов, хотя на графике это одно облако. Сумма квадратов расстояний до центров при этом уменьшилась мало, с 67,8 до 57,9. Оптимальное число кластеров — 3.

\subsection{Метод локтя}

Для каждого числа кластеров $k$ от 1 до 9 вычислена сумма квадратов расстояний от объектов до центров их кластеров по~\formref{inertia}. В модели KMeans она хранится в атрибуте inertia\_.

\suaieq[inertia]{J(C) = \sum_{k=1}^{K} \sum_{i \in C_k} \lVert x_i - \mu_k \rVert^2}
K | число кластеров
C_k | множество объектов кластера $k$
x_i | вектор признаков объекта $i$
\mu_k | центр кластера $k$

С ростом $k$ эта сумма всегда убывает, поэтому ищется не минимум, а «локоть» — точка, после которой убывание резко замедляется. График показан на~\figref{p1_elbow}, значения и их снижение относительно предыдущего $k$ — на~\figref{p1_elbow_table}.

\suaiimg[0.75\textwidth]{p1_elbow}{График по методу локтя для сгенерированных данных}

\suaiimg[0.52\textwidth]{p1_elbow_table}{Значения $J(C_k)$ и их снижение}

При переходе от одного кластера к двум сумма уменьшается на 46,2~\%, от двух к трём — ещё на 79,0~\%. Дальше каждый новый кластер снижает её только на 12–15~\%. Локоть находится в точке $k = 3$.

\subsection{Выводы по части 1}

Метод k-means для трёх кластеров после стандартизации точно восстановил сгенерированные группы: в каждом кластере по 100 объектов, ARI~= 1, центры кластеров в исходных единицах близки к заложенным в функцию. Коэффициент силуэта максимален при трёх кластерах: 0,720 против 0,499 для двух и 0,585 для четырёх. При двух кластерах алгоритм объединяет две группы, при четырёх делит одну группу пополам. Метод локтя указывает на то же число: после $k = 3$ сумма квадратов расстояний почти перестаёт убывать.

\section{Кластеризация набора данных о пациентах}

\subsection{Загрузка данных}

Набор данных medical\_diagnosis.csv загружен методом read\_csv. Результат метода info показан на~\figref{data_info}.

\suaiimg[\textwidth]{data_info}{Общая информация о наборе данных medical\_diagnosis.csv}

В таблице 3031 строка и 14 столбцов, пропусков нет. Числовых столбцов 9, включая целевой столбец Disease, категориальных 5. Типы столбцов соответствуют их содержанию. Одна строка полностью повторяет другую.

\subsection{Предварительная обработка данных}

Описательные статистики числовых столбцов показаны на~\figref{describe}, значения категориальных столбцов и число строк с каждым значением — на~\figref{categorical}.

\suaiimg[\textwidth]{describe}{Описательные статистики числовых столбцов}

\suaiimg[0.49\textwidth]{categorical}{Значения категориальных столбцов}

У столбца cholesterol\_mg\_dl максимум 658,2 почти втрое больше третьего квартиля (232). У столбца CRP\_mg\_L максимум 43,45 при медиане 0,52, а стандартное отклонение (3,16) больше среднего (1,38): распределение сильно скошено вправо. У остальных признаков среднее близко к медиане, а минимум и максимум правдоподобны. В категориальных столбцах встречаются только значения, указанные в описании набора. В кластеризации категориальные признаки не участвуют: k-means и метод Уорда используют евклидово расстояние, а оно определено для числовых признаков.

Для проверки двух столбцов с подозрительно большими максимумами построена диаграмма рассеяния холестерина и C-реактивного белка (рисунок~\ref{fig:anomalies}).

\suaiimg[0.75\textwidth]{anomalies}{Холестерин и C-реактивный белок до обработки}

Основная масса пациентов лежит в левом нижнем углу. Справа вверху отдельно расположена группа точек, у которых одновременно холестерин больше 300 и C-реактивный белок больше 20. Результаты проверки этой группы показаны на~\figref{anomalies_check}.

\suaiimg[\textwidth]{anomalies_check}{Проверка группы строк с завышенными холестерином и C-реактивным белком}

В группе 31 строка, и все они записаны подряд в самом конце файла, под номерами с 3000 по 3030. У 26 из них холестерин дробный, а во всех остальных 3000 строках он целый. Медиана холестерина в группе равна 487,2 и превышает максимум среди остальных строк (364), медиана C-реактивного белка — 24,48 против 0,51. Остальные шесть признаков лежат в тех же пределах, что и у остальных пациентов. Обе повторяющиеся строки (3008 и 3024) также входят в эту группу. Положение в файле, другой формат значений и одновременное завышение двух признаков показывают, что это искажённые записи, а не пациенты с редкими анализами.

Для кластеризации такие строки вредны по двум причинам. Центр кластера является средним, и далёкие точки смещают его к себе. Кроме того, они увеличивают стандартное отклонение признака, и после стандартизации различия между обычными пациентами по этому признаку сжимаются. Для проверки метод k-means для трёх кластеров применён к данным вместе с этими строками (рисунок~\ref{fig:anomalies_kmeans}).

\suaiimg[\textwidth]{anomalies_kmeans}{Кластеризация данных вместе с аномальными строками}

Алгоритм выделил в отдельный кластер 31 аномальную строку и ещё три строки с самым высоким C-реактивным белком, а остальных 2997 пациентов разделил пополам. Аномальные строки удалены (рисунок~\ref{fig:cleaned}).

\suaiimg[\textwidth]{cleaned}{Удаление аномальных строк}

Осталось 3000 строк, дубликат исчез вместе с аномальными строками. Для оставшихся строк подсчитано число выбросов по правилу полутора межквартильных размахов и число строк, в которых значение признака равно его минимуму (рисунок~\ref{fig:outliers}).

\suaiimg[0.78\textwidth]{outliers}{Выбросы и минимальные значения числовых признаков}

У семи признаков из восьми выбросов от 0 до 16, асимметрия близка к нулю: это края обычных распределений. У CRP\_mg\_L 220 выбросов и асимметрия 4,54: у половины пациентов значение не больше 0,51, у немногих оно доходит до 26,2. Повышенный C-реактивный белок является признаком воспаления, то есть сведением о состоянии пациента, а не ошибкой. Минимальные значения повторяются: у 857 пациентов CRP\_mg\_L равен 0,1, у 287 symptom\_score равен 0, у 103 глюкоза равна 56. Вероятно, это нижние границы шкал, которыми заменены все меньшие значения. Эти строки оставлены.

\subsection{Целевая переменная и матрица диаграмм рассеяния}

Целевая переменная — Disease: 1 означает, что пациент болеет, 0 — не болеет. В модель она не включена, так как кластеризация относится к обучению без учителя и группы должны быть найдены только по признакам. Целевая переменная используется далее для сравнения кластеров с известными классами. Матрица признаков составлена из восьми числовых столбцов (рисунок~\ref{fig:target}).

\suaiimg[\textwidth]{target}{Выделение целевой переменной и матрицы признаков}

Классы сильно не сбалансированы: больных 48 из 3000, то есть 1,6~\%. Матрица диаграмм рассеяния числовых признаков построена функцией pairplot (рисунок~\ref{fig:pairplot}). Красным цветом выделены больные, голубым — здоровые. Показана нижняя половина матрицы, так как верхняя её зеркально повторяет.

\suaiimg[\textwidth]{pairplot}{Матрица диаграмм рассеяния с выделением целевой переменной}

Ни на одной диаграмме нет отдельных сгустков точек: пациенты образуют одно сплошное облако, явных кластеров в данных не видно. Облака не вытянуты ни в каком направлении, то есть признаки между собой не связаны. Красные точки не образуют отдельной группы и разбросаны внутри голубого облака, но на диагонали кривая распределения больных по каждому признаку сдвинута вправо. Заметнее всего отличие по CRP\_mg\_L: красные точки лежат полосой выше основной массы голубых. Средние значения признаков у здоровых и больных приведены на~\figref{disease_means}.

\suaiimg[\textwidth]{disease_means}{Средние значения признаков у здоровых и больных}

У больных все восемь признаков в среднем выше: возраст 65,7 года против 52,9, индекс массы тела 31,6 против 26,4, давление 139 против 122, глюкоза 116 против 94, C-реактивный белок 3,90 против 1,07. При этом ни один признак не разделяет классы сам по себе. У всех больных C-реактивный белок не ниже 2,16, но такое же или большее значение имеют 15,4~\% всех пациентов.

\subsection{Стандартизация}

Признаки измеряются в разных единицах: стандартное отклонение холестерина равно 45,01, а symptom\_score — 1,36. Без стандартизации расстояние между пациентами почти полностью определялось бы признаками с большим разбросом. Все восемь признаков приведены к среднему 0 и стандартному отклонению 1 (рисунок~\ref{fig:scale}).

\suaiimg[\textwidth]{scale}{Стандартизация числовых признаков}

\subsection{Подбор числа кластеров для метода k-means}

Для числа кластеров $k$ от 1 до 10 обучены модели k-means и вычислены сумма квадратов расстояний до центров $J(C_k)$ и коэффициент силуэта. Для одного кластера коэффициент силуэта не определён. Значения приведены на~\figref{k_selection_table}, графики — на~\figref{k_selection}.

\suaiimg[0.72\textwidth]{k_selection_table}{Показатели моделей k-means для разного числа кластеров}

\suaiimg[\textwidth]{k_selection}{Метод локтя и коэффициент силуэта для набора данных о пациентах}

Выраженного локтя на графике нет: сумма убывает плавно, каждый новый кластер снижает её на всё меньшую долю — 9,5~\%, 7,3~\%, 6,4~\%, 5,5~\% и далее. В части 1 переход к трём кластерам снижал сумму на 79~\%, здесь наибольший шаг равен 9,5~\%. Коэффициент силуэта максимален при $k = 3$ и равен 0,105. При $k = 2$ он равен 0,100, при остальных $k$ — от 0,092 до 0,096. Оптимальное число кластеров — 3.

Само значение коэффициента мало: близость к нулю означает перекрывающиеся кластеры. Это согласуется с матрицей диаграмм рассеяния, на которой пациенты образуют одно облако. При $k = 3$ кроме двух больших кластеров появляется маленький, из 185 пациентов. Похожий маленький кластер, от 117 до 164 пациентов, сохраняется и при всех больших $k$.

\subsection{Анализ кластеров метода k-means}

Итоговая модель обучена для трёх кластеров, метки кластеров добавлены в таблицу столбцом cluster\_kmeans (рисунок~\ref{fig:kmeans_fit}).

\suaiimg[\textwidth]{kmeans_fit}{Обучение итоговой модели k-means и размеры кластеров}

Алгоритм сошёлся за 31 итерацию. Получены два больших кластера, 1452 и 1363 пациента, и один маленький — 185 пациентов. Для описания кластеров написаны три функции (листинг~\ref{lst:helpers}). Функция cluster\_means вычисляет средние значения признаков и долю больных в каждом кластере и по всему набору. Функция clustering\_metrics вычисляет метрики качества, функция disease\_by\_cluster — число здоровых и больных в каждом кластере.

\suaicode[Python, helpers]{Функции для описания кластеров и оценки качества}
    def cluster_means(column):
        means = medical.groupby(column)[NUMERIC + [TARGET]].mean().T
        means['Все пациенты'] = medical[NUMERIC + [TARGET]].mean()
        return means.round(2)


    def clustering_metrics(labels):
        return {'Силуэт': metrics.silhouette_score(X_scaled, labels),
                'ARI': metrics.adjusted_rand_score(y, labels),
                'AMI': metrics.adjusted_mutual_info_score(y, labels),
                'FMI': metrics.fowlkes_mallows_score(y, labels),
                'Homogeneity': metrics.homogeneity_score(y, labels),
                'Completeness': metrics.completeness_score(y, labels),
                'V-measure': metrics.v_measure_score(y, labels)}


    def disease_by_cluster(column):
        table = pd.crosstab(medical[column], medical[TARGET])
        share = 100 * table[1] / table.sum(axis=1)
        table['Доля больных, %'] = share.round(1)
        return table

Средние значения признаков в кластерах в исходных единицах приведены на~\figref{kmeans_means}. Последняя строка, среднее значение Disease, равна доле больных.

\suaiimg[0.48\textwidth]{kmeans_means}{Средние значения признаков в кластерах k-means}

Центры кластеров в стандартизованных признаках показывают, на сколько стандартных отклонений кластер отличается от среднего пациента (рисунок~\ref{fig:kmeans_centers}).

\suaiimg[\textwidth]{kmeans_centers}{Центры кластеров k-means в стандартизованных признаках}

Кластеры различаются по двум признакам. По возрасту центры кластеров 1 и 2 отстоят от среднего на $-0{,}83$ и $+0{,}87$ стандартного отклонения, по C-реактивному белку центр кластера 0 — на $+2{,}93$. По остальным шести признакам отклонения не больше 0,20. Границы кластеров по этим двум признакам приведены на~\figref{kmeans_bounds}, диаграмма рассеяния в осях возраст — C-реактивный белок — на~\figref{kmeans_plot}.

\suaiimg[\textwidth]{kmeans_bounds}{Минимум, медиана и максимум возраста и C-реактивного белка в кластерах k-means}

\suaiimg[0.8\textwidth]{kmeans_plot}{Кластеры k-means в осях возраст — C-реактивный белок}

Кластер 0 составляют пациенты с повышенным C-реактивным белком: 185 человек, или 6,2~\% набора. Значение у всех не ниже 3,51, в среднем 6,20 — в 5,5 раза выше среднего по набору. Возраст любой, от 19 до 89 лет, глюкоза немного выше средней (98,3 против 94,4).

Кластер 1 составляют молодые пациенты: 1452 человека, или 48,4~\%. Средний возраст 35,7 года, самому старшему 65 лет. C-реактивный белок низкий: в среднем 0,79 и не выше 4,00.

Кластер 2 составляют пожилые пациенты: 1363 человека, или 45,4~\%. Средний возраст 71,1 года, самому младшему 43 года. C-реактивный белок такой же низкий, в среднем 0,78.

На диаграмме видно, что кластер 0 отделён от остальных по уровню C-реактивного белка, граница проходит около 3,5–4~мг/л. Кластеры 1 и 2 разделены по возрасту, граница проходит около 53 лет. Пустого пространства между кластерами нет: границы проведены внутри сплошного облака точек. По давлению, холестерину, индексу массы тела и лейкоцитам кластеры не различаются.

\subsection{Метрики качества}

Коэффициент силуэта является внутренней метрикой: для него нужны только признаки и метки кластеров. Внешние метрики сравнивают кластеры с известными классами Disease. Метрика ARI вычисляется по~\formref{ari}:
\suaieq[ari]{ARI = \frac{RI - E[RI]}{\max(RI) - E[RI]}}
RI | индекс Рэнда — доля пар объектов, по которым два разбиения согласны
E[RI] | ожидаемое значение индекса Рэнда при случайном разбиении

ARI равен 1 при совпадении разбиений и близок к 0 при случайном согласии. Метрика AMI устроена так же, но основана на взаимной информации двух разбиений. Индекс Фаулкса — Мэллоуса FMI вычисляется по~\formref{fmi} и принимает значения от 0 до 1:
\suaieq[fmi]{FMI = \frac{TP}{\sqrt{(TP + FP)(TP + FN)}}}
TP | число пар объектов одного класса, попавших в один кластер
FP | число пар объектов разных классов, попавших в один кластер
FN | число пар объектов одного класса, попавших в разные кластеры

Однородность (homogeneity) равна 1, если каждый кластер состоит из объектов одного класса, и вычисляется по~\formref{homogeneity}. Полнота (completeness) равна 1, если все объекты одного класса попали в один кластер, и вычисляется так же с перестановкой классов и кластеров. V-мера является их гармоническим средним и вычисляется по~\formref{vmeasure}.

\suaieq[homogeneity]{h = 1 - \frac{H(C \mid K)}{H(C)}}
H(C \mid K) | условная энтропия классов при известном разбиении на кластеры
H(C) | энтропия классов

\suaieq[vmeasure]{V = \frac{2hc}{h + c}}
h | однородность
c | полнота

Метрики разбиения k-means приведены на~\figref{kmeans_metrics}, распределение здоровых и больных по кластерам — на~\figref{kmeans_disease}.

\suaiimg[0.6\textwidth]{kmeans_metrics}{Метрики качества кластеризации методом k-means}

\suaiimg[0.35\textwidth]{kmeans_disease}{Здоровые и больные в кластерах k-means}

Коэффициент силуэта равен 0,105: кластеры перекрываются. ARI и AMI равны 0,012, то есть совпадение кластеров с классами Disease почти не отличается от случайного. Однородность равна 0,076, полнота 0,007, V-мера 0,013: в каждом кластере преобладают здоровые, а больные распределены по всем трём кластерам — 15, 9 и 24 человека. Значение FMI, равное 0,661, является следствием дисбаланса классов. Здоровых 98,4~\%, поэтому почти любая пара пациентов принадлежит одному классу, и большое значение получается при любом разбиении на несколько крупных кластеров.

Кластеры не воспроизводят деление на больных и здоровых. Алгоритм уменьшает разброс внутри кластеров, а больные составляют 1,6~\% набора и на суммарный разброс почти не влияют. При этом кластер 0 содержателен: доля больных в нём равна 8,1~\%, что в пять раз выше, чем по всему набору. В кластере молодых пациентов доля больных наименьшая, 0,6~\%, в кластере пожилых — 1,8~\%.

\subsection{Иерархическая агломеративная кластеризация}

Агломеративный алгоритм сначала считает каждый объект отдельным кластером и на каждом шаге объединяет два ближайших кластера, пока не останется один. Использован метод Уорда: объединяются два кластера, слияние которых меньше всего увеличивает внутрикластерную дисперсию. Он уменьшает ту же величину, что и k-means, поэтому результаты двух методов сравнимы. Расстояние между объектами евклидово.

История объединений получена функцией linkage библиотеки scipy, по ней функцией dendrogram построена дендрограмма (рисунок~\ref{fig:dendrogram}). По горизонтали расположены 3000 пациентов, по вертикали — расстояние, на котором кластеры объединились.

\suaiimg[\textwidth]{dendrogram}{Дендрограмма, метод Уорда}

Чем длиннее вертикальные ветви между двумя соседними объединениями, тем сильнее различаются кластеры, которые они разделяют. Расстояния последних десяти объединений приведены на~\figref{merges}.

\suaiimg[0.73\textwidth]{merges}{Расстояния последних десяти объединений}

Наибольший скачок расстояния равен 7,89. Четыре кластера объединяются в три на расстоянии 41,06, а три в два — только на расстоянии 48,95. Между этими уровнями дерево можно разрезать, не задев ни одного объединения: красная пунктирная линия на уровне 45 пересекает три вертикальные ветви. Оптимальное число кластеров по дендрограмме — 3. Следующие по величине скачки, 6,29 и 6,04, соответствуют восьми и двум кластерам.

Для разрезов дерева на 2–6 кластеров вычислен коэффициент силуэта (рисунок~\ref{fig:ward_selection}).

\suaiimg[0.42\textwidth]{ward_selection}{Коэффициент силуэта для разрезов дерева}

По коэффициенту силуэта лучше два кластера: 0,270. Но такое разбиение почти ничего не сообщает о пациентах: от 2871 человека отделены 129 с самым высоким C-реактивным белком. Коэффициент высок, потому что для 96~\% объектов единственный чужой кластер — эта далёкая маленькая группа. При трёх кластерах большой кластер делится на два, и коэффициент падает до 0,060. Выбраны три кластера: на них указывает наибольший скачок на дендрограмме, и результат можно напрямую сравнить с k-means.

Модель AgglomerativeClustering с методом Уорда обучена для трёх кластеров, метки добавлены в таблицу столбцом cluster\_agglo (рисунок~\ref{fig:agglo_fit}). Размеры кластеров совпадают с разрезом дерева на три кластера.

\suaiimg[\textwidth]{agglo_fit}{Агломеративная кластеризация и размеры кластеров}

Средние значения признаков в кластерах приведены на~\figref{agglo_means}, границы кластеров по возрасту и C-реактивному белку — на~\figref{agglo_bounds}. Номера кластеров у двух методов между собой не связаны.

\suaiimg[\textwidth]{agglo_means}{Средние значения признаков в кластерах агломеративной кластеризации}

\suaiimg[\textwidth]{agglo_bounds}{Минимум, медиана и максимум возраста и C-реактивного белка в кластерах агломеративной кластеризации}

Кластер 1 составляют пациенты с повышенным C-реактивным белком: 129 человек, или 4,3~\%, среднее значение 6,87, наименьшее 2,28. Это аналог кластера 0 метода k-means, но меньшего размера. В кластере 0 находятся 1609 человек (53,6~\%): пациенты моложе среднего (43,6 года), с глюкозой ниже средней (88,2 против 94,4) и более выраженными симптомами (2,27 против 2,04). В кластере 2 находятся 1262 человека (42,1~\%): пациенты старше среднего (64,5 года), с глюкозой выше средней (101,7) и менее выраженными симптомами (1,76). В отличие от k-means, кластеры 0 и 2 не разделены по возрасту: в обоих есть пациенты от 18 до 89 лет. Они различаются сразу по трём признакам, но по каждому слабо.

Метрики качества обоих методов приведены на~\figref{metrics_both}, распределение здоровых и больных по кластерам — на~\figref{agglo_disease}.

\suaiimg[\textwidth]{metrics_both}{Метрики качества двух методов кластеризации}

\suaiimg[\textwidth]{agglo_disease}{Здоровые и больные в кластерах агломеративной кластеризации}

Коэффициент силуэта агломеративной кластеризации равен 0,060 и почти вдвое ниже, чем у k-means. Внешние метрики так же малы: ARI~= 0,009, AMI~= 0,011, V-мера 0,012. Деление на больных и здоровых этот метод тоже не воспроизводит. В кластере с повышенным C-реактивным белком доля больных равна 4,7~\% (6 из 129), а у k-means в таком же кластере — 8,1~\% (15 из 185).

\subsection{Сравнение методов}

Для сравнения двух разбиений между собой вычислен ARI и построена таблица сопряжённости меток (рисунок~\ref{fig:methods_crosstab}).

\suaiimg[\textwidth]{methods_crosstab}{Сравнение разбиений k-means и агломеративной кластеризации}

ARI между разбиениями равен 0,271: методы согласны лишь отчасти. Группу с повышенным C-реактивным белком находят оба метода: из 129 пациентов кластера 1 по методу Уорда 118 входят в кластер 0 метода k-means. Два больших кластера проведены по-разному: из кластера молодых пациентов k-means 304 человека попали в старший кластер метода Уорда, а из кластера пожилых 419 человек — в младший. Оба разбиения показаны на~\figref{methods_plot}.

\suaiimg[\textwidth]{methods_plot}{Кластеры двух методов в осях возраст — C-реактивный белок}

На левой диаграмме граница по возрасту чёткая, на правой синие и зелёные точки перемешаны. Причина расхождения заключается в самих данных: явных кластеров в них нет, и границу внутри сплошного облака каждый алгоритм проводит по-своему. Итоговые показатели двух методов сведены в~\tabref{methods}.

\suaitable[methods]{Сравнение методов кластеризации}
Показатель | k-means | Агломеративная кластеризация
Число кластеров | 3 | 3
Размеры кластеров | 1452, 1363, 185 | 1609, 1262, 129
Коэффициент силуэта | 0,105 | 0,060
ARI с классами Disease | 0,012 | 0,009
Доля больных в кластере с повышенным C-реактивным белком, \% | 8,1 | 4,7

Целесообразнее использовать метод k-means. Коэффициент силуэта у него почти вдвое выше. Его кластеры проще описать: молодые пациенты, пожилые пациенты и пациенты с повышенным C-реактивным белком, тогда как кластеры метода Уорда различаются по нескольким признакам сразу и по каждому слабо. Метод k-means на каждой итерации заново распределяет объекты по кластерам, а агломеративный алгоритм не пересматривает объединение, сделанное на раннем шаге. Кроме того, k-means лучше подходит для больших выборок: агломеративному алгоритму нужны попарные расстояния между всеми объектами, а их число растёт как квадрат размера выборки.

Иерархический метод полезен на другом этапе: дендрограмма показывает структуру данных и помогает выбрать число кластеров, не задавая его заранее. В этой работе она подтвердила выбор трёх кластеров.

\suaiconclusion

В ходе работы изучены два метода кластеризации, k-means и иерархическая агломеративная кластеризация, и метрики качества кластеризации.

В первой части на 300 сгенерированных объектах трёх групп метод k-means после стандартизации точно восстановил группы (ARI~= 1). Коэффициент силуэта максимален при трёх кластерах (0,720 против 0,499 для двух и 0,585 для четырёх), метод локтя указывает на то же число.

Во второй части исследован набор medical\_diagnosis.csv: 3031 запись о пациентах с восемью числовыми и пятью категориальными признаками и целевым классом Disease. Пропусков нет. Удалена 31 аномальная строка с одновременно завышенными холестерином и C-реактивным белком, включая единственный дубликат. Осталось 3000 строк, из них 48 больных (1,6~\%). Восемь числовых признаков стандартизованы. По матрице диаграмм рассеяния пациенты образуют одно сплошное облако.

Для k-means коэффициент силуэта максимален при трёх кластерах и равен 0,105. Получены кластеры молодых пациентов (1452 человека, средний возраст 35,7 года), пожилых (1363 человека, 71,1 года) и пациентов с повышенным C-реактивным белком (185 человек). В последнем доля больных равна 8,1~\% и в пять раз превышает среднюю. Агломеративная кластеризация методом Уорда по дендрограмме также даёт три кластера, но с коэффициентом силуэта 0,060. Внешние метрики обоих методов близки к нулю (ARI 0,012 и 0,009): кластеры не совпадают с делением на больных и здоровых.

Целесообразнее использовать k-means: у него выше коэффициент силуэта, а кластеры проще интерпретировать. Качество кластеризации определяется самими данными: отделённые группы первой части k-means находит безошибочно, а в данных о пациентах таких групп нет, и кластеры являются условными сегментами внутри сплошного облака. Диагноз кластеризация не восстанавливает, но выделяет группу риска — пациентов с повышенным C-реактивным белком.

\suaisources
Учебно-методические указания к выполнению лабораторной работы №4 по дисциплине «Введение в анализ данных».

\end{document}
